\chapter*{List of Symbols and Abbreviations}
\addcontentsline{toc}{chapter}{List of Symbols and Abbreviations}

\small
\begin{tabularx}{0.9\textwidth}{@{}p{0.22\textwidth}@{\hspace{0.4em}:\hspace{0.4em}}X@{}}
    \textbf{Symbol or} \newline \textbf{Abbreviation} & \textbf{Explanation}\\
    
    %% Edit below in the format
    %% XYZ &:& EXPLANATION\\
    %% where XYZ is the acronym or symbol
    %% and EXPLANATION is the explanation of it
    %% make sure not to forget &:& between them, and \\ at the end of EXPLANATION
    
    AES & Advanced Encryption Standard\\
    AS608 & Optical fingerprint sensor module used in the project\\
    CBC & Cipher Block Chaining\\
    CLI & Command-Line Interface\\
    CTR & Counter Mode\\
    dm-crypt & Linux kernel device-mapper encryption target\\
    ECB & Electronic Codebook\\
    ext4 & Fourth extended filesystem used inside the unlocked vault\\
    FAT32 & File Allocation Table 32-bit filesystem\\
    FUSE & Filesystem in Userspace\\
    GUI & Graphical User Interface\\
    HCI & Human-Computer Interaction\\
    IV & Initialization Vector\\
    LUKS & Linux Unified Key Setup\\
    LUKS2 & Version 2 of Linux Unified Key Setup\\
    NTFS & New Technology File System\\
    PC & Personal Computer\\
    QML & Qt Modeling Language\\
    Qt & Cross-platform application development framework\\
    RAM & Random Access Memory\\
    STM32 & ARM Cortex-M based microcontroller family used in the project\\
    TPM & Trusted Platform Module\\
    UART & Universal Asynchronous Receiver/Transmitter\\
    USB & Universal Serial Bus\\
    XTS & XEX-based Tweaked-codebook mode with ciphertext stealing\\
    
    %% Until here
\end{tabularx}
\normalsize

\clearpage
